\scanpage{79}
Donc $\Sigma$ se réconstitue à partir de $S$ comme l'\deleted{\unsure{ens. des|ens. des ob.}} sous-catégorie (str.) pleine des objets ``de présentation finie'' !

\[
\begin{aligned}
S:\qquad&\Sigma=S_{pF}\subset S\\
&R=\Sigma^\circ\\
&T:C\longmapsto\underline{\mathrm{Hom}}_{\mathrm{lex}}(S_{pF}^{\,\circ},\unsure{C|e})
\end{aligned}
\]

\[
\begin{aligned}
T:\qquad&S=T(\mathrm{Ens})\\
&\Sigma=S_{pF}=T(\mathrm{Ens})_{Fp}\\
&R=\Sigma^\circ
\end{aligned}
\]

\uline{Corollaire} Pour qu'une \deleted{\unsure{th.|théorie}} $\Delta_1$-théorie \inserted{$T$} (pour $\varprojlim$ quelc) provienne d'une théorie définissable par $\varprojlim$ finie, il f. et s. que $S=S_T=T(\mathrm{Ens})$ soit \deleted{\unsure{engendré}} strictement engendrée par \deleted{ses} \inserted{sa sous-catégorie pleine $\Sigma$ \unsure{de ses|des}} objets de \unsure{prés.|Pr.} finie, ou encore que $S\longrightarrow\deleted{\unsure{\mathrm{Ens}}}\widehat{\Sigma}$ soit pleinement fidèle. (Alors $\Sigma^\circ$ est la catégorie classifiante \deleted{pour} la théorie sur $\Delta_0$ qui donne naissance à $T$).

\uline{Corollaire} Le 2-foncteur de restriction à $\Delta_1$ des théories algébriques sur $\Delta_0$ est 1-fidèle, i.e. pour \deleted{\unsure{la}} deux théories $T,T'$ sur $\Delta_0$, on a
\[
\underline{\mathrm{Hom}}_{\mathrm{Th}_{\deleted{\unsure{\Delta_0}}\Delta_0}}(T,T')\longrightarrow\underline{\mathrm{Hom}}_{\mathrm{Th}_{\Delta_1}}(T_1,T'_1)\quad\text{pl. fidèle.}
\]
L'image essentielle est formée des foncteurs $T_1-T'_1$ qui sont tels que
\[
S=T(\mathrm{Ens})\inserted{=T_1(\mathrm{Ens})}\longrightarrow S'=T'(\mathrm{Ens})=T'_1(\mathrm{Ens})
\]
transforme objets de P.F. en objets de P.F. [ On pourrait les appeler les morphismes de présent. finie, ou cohérents, de théories \ldots ].

% vertical emphasis stroke beside the question
\uline{Question} Considérons $\Delta$ définie par $d'=$ ens des diagrammes finis \uline{discrets}, $d''=\varnothing$, pas de condition d'exactitude. Soit $R$ une petite cat $\in\mathrm{Ob}\,\Delta$, $\Sigma=R^\circ\subset S=\unsure{T_R|\mathrm{Th}_R}(\mathrm{Ens})=\underline{\mathrm{Hom}}_{\mathrm{prod}}(R,\mathrm{Ens})$. Peut-il se faire que le $=$ $S$ soit défini par des sous-catégories \deleted{essentielles}\inserted{str. pleines} différentes ? \deleted{Oui, comme on voit avec} Comment trouver \inserted{toutes} les $\Sigma$ qui marchent, i.e. tels que $S\cong\rho_{\Delta,\Delta_1}(\Sigma)$ ? on aura évidem\textsuperscript{mt} $\Sigma\subset S_{pF}=\Sigma_0$ et $\unsure{\Sigma_{pF}|\Sigma_0|S_{pF}}\cong\rho_{\Delta,\Delta_0}(\Sigma)$ (condition néc. et suff.) Donc la question se repose en termes d'un $\Sigma_0\in\mathrm{Ob}\,\Delta_0$ et d'une \deleted{famille \unsure{d'obj.|génératrice} dans} \inserted{sous-catégorie pleine} $\Sigma$ de $\Sigma_0$ stable par sommes finies\ldots
% archive sheet number: 78
